\scanpage{524}
\mspage{435}
% B01
% slanted
\margin{\uline{faux} ?}

% B02
\uline{Proposition} 4.15.6. Soit $f:E\unsure{\relbar|\to} C$ un foncteur de $\mathcal U$-catégories. On suppose

% B03
a) $f$ est fibrant ou cofibrant, \deleted{\unsure{et les}}

% B04
% slanted
\margin{Le cas fibrant n'est pas prouvé, peut-être \uuline{faux}. -- et de $=$ pour le cas cofibrant !}

% B05
b) les catégories fibres $E_s$ ($s\in\mathrm{Ob}\,C$) sont accessibles, et pour tte flèche $\lambda:s\relbar t$ dans $C$, le foncteur chgt de base $\lambda^*:E_t\relbar E_s$ (resp. cochgt. de base $\lambda_*:E_s\relbar E_t$) est accessible.

% B06
c) $C$ est ess. petite et karoubienne

Alors \deleted{\unsure{$f$}}\unsure{$E$} est accessible, et $f$ est accessible.

% B07
\uline{Dém} On a déjà vu que les conditions a) et b) impliquent $H'$, i.e. $H(\unsure{1+2|1,2}+3)$. D'ailleurs, si $C'$ est une \inserted{\deleted{\unsure{petite}}} catégorie \deleted{\unsure{petite}} et $C'\xrightarrow{\approx}C$ une équivalence, alors $E'\overset{\mathrm{déf}}{=}E\times_C C'\relbar E$ est une équivalence, donc la conclusion voulue pour $E,f$ \deleted{\unsure{revient}} équivaut à la $=$ conclusion pour $\unsure{E'|C'},f':E'\relbar C'$, \deleted{\unsure{OPS}} Donc OPS que $C$ est petite. (et $=$ réduite, si on y tient).
% B08
En vertu de 4.15.\unsure{2|1} a), on voit donc que les foncteurs d'inclusion $E_s\longrightarrow E$ sont accessibles (et $=\pi_1$-accessibles dès que $\pi_1\geq\mathop{\mathrm{Sup}}\limits_{\lambda\in\mathrm{Fl}\,C}\pi_\lambda$). On voit d'ailleurs

% archive sheet number: 523
